\documentclass{article}
\usepackage{amsmath,amssymb}
\usepackage{xurl}
\usepackage[hidelinks]{hyperref}
% Shared commands for notes in docs/paper-gaps/.

\DeclareUnicodeCharacter{2080}{\ifmmode{}_{0}\else\textsubscript{0}\fi}
\DeclareUnicodeCharacter{2081}{\ifmmode{}_{1}\else\textsubscript{1}\fi}
\DeclareUnicodeCharacter{2082}{\ifmmode{}_{2}\else\textsubscript{2}\fi}
\DeclareUnicodeCharacter{2099}{\ifmmode{}_{n}\else\textsubscript{n}\fi}

\newcommand{\Lean}{\textsc{Lean}}
\ExplSyntaxOn
\NewDocumentCommand{\leanid}{m}{
  \ifmmode
    \textsf{\detokenize{#1}}
  \else
    \group_begin:
    \urlstyle{sf}
    \tl_set:Nn \l_tmpa_tl {#1}
    \tl_replace_all:Nnn \l_tmpa_tl {\_} {_}
    \exp_args:NV \path \l_tmpa_tl
    \group_end:
  \fi
}
\ExplSyntaxOff
\providecommand{\tr}{\operatorname{tr}}
\newcommand{\tnleanIssueBase}{https://github.com/LionSR/TNLean/issues/}
\newcommand{\tnleanPrBase}{https://github.com/LionSR/TNLean/pull/}
\newcommand{\tnleanDocBase}{https://sirui-lu.com/TNLean/docs/}
\newcommand{\ghissue}[1]{\href{\tnleanIssueBase#1}{GitHub issue \##1}}
\newcommand{\ghpr}[1]{\href{\tnleanPrBase#1}{PR \##1}}
\newcommand{\doclink}[1]{\href{\tnleanDocBase#1.html}{\path{#1}}}

% Dirac notation and the identity operator, as in blueprint/src/macros/common.tex.
% \providecommand keeps note-local definitions authoritative.
\usepackage{dsfont}
\providecommand{\ket}[1]{|#1\rangle}
\providecommand{\bra}[1]{\langle #1|}
\providecommand{\braket}[2]{\langle #1 | #2 \rangle}
\providecommand{\ketbra}[2]{|#1\rangle\!\langle #2|}
\providecommand{\Id}{\mathds{1}}
\providecommand{\one}{\mathds{1}}
\providecommand{\C}{\mathbb{C}}
\providecommand{\MN}[1]{M_{#1}(\C)}
\providecommand{\MD}{M_D(\C)}

% Verdict marker: \gapnote{<kind>}{<status>}. Kind is the policy
% classification (clarification, local-correction, scope-restriction,
% unfaithful, false-source, open-gap); status is open, wip (elimination
% underway), resolved, or historical. Machine-read by the site index;
% prints a verdict line.
\newcommand{\gapnote}[2]{\par\noindent\textbf{Verdict:} #1 (#2).\par}

\title{Inhomogeneous Long-Range Preparation: Coherent Pair Products}
\author{}
\date{2026-11-15}
\begin{document}
\maketitle
\gapnote{scope-restriction}{open}

\section{The source remark}
The concluding sentence of ``Long-range MPS using measurements'' in
arXiv:2307.01696v2, page~5, immediately after equation~(19), states:
``Our construction generalizes to inhomogeneous long-range correlated MPS
exactly as in the short-range case.'' The issue cites
\texttt{References/2307.01696/main.tex:913}; that source directory is absent
from the repository snapshot, so the sentence and its context were checked
against the published arXiv PDF and HTML.

The preceding short-range paragraph defines a sequence by approximation of
the blocked positive-part state by site-dependent nearest-neighbour pair
products at a logarithmic blocking scale, with error tending to zero. The
long-range paragraph supplies a coherent sum of pair products in (19),
with orthonormal label states at each pair window. Its final sentence does
not separately specify an inhomogeneous long-range class, a sector
decomposition for it, or an error rate uniform over the sequence.

\section{The explicit conditional construction}
The new result in
\doclink{TNLean/MPS/Preparation/InhomogeneousLongRange} assumes
\begin{align*}
    \Omega'=\sum_{j=1}^b\alpha_j\bigotimes_{k=1}^M\omega_j^k,
    \qquad \sum_j|\alpha_j|^2=1,
    \qquad \langle\omega_i^k,\omega_j^k\rangle=\delta_{ij}.
\end{align*}
For physical dimension $d\ge2$, common bond dimension $D\ge1$, and blocks
of lengths $3(b+D+2)\le\ell_k\le L$, the written construction chooses
isometric extensions $X_k\Pi_k=V_k$ of the blocked polar factors before
choosing the pairs or the amplitudes. It prepares
$\psi=(\bigotimes_kX_k)\Omega'$ with measurements followed by a circuit
in depth $CL$, with $C$ uniform in the ring length, number of blocks,
tensors, pairs and amplitudes. No injectivity is required. The overlap
with the normalized target is exactly the overlap of $\Omega'$ with the
normalized positive-part state. Both tensors and pairs may vary along
the ring, and the circuit accounts for every physical site.

This is a construction with specified input assumptions, rather than a
formal definition of all inhomogeneous long-range correlated MPS. Its
blueprint statement has its own label and does not certify the unrestricted
source remark.

\section{Remaining obligations}
An explicit one-partition condition is now written as
\leanid{MPSPreparation.IsCoherentPairApproximable} in
\doclink{TNLean/MPS/Preparation/InhomogeneousLongRangeRate}. It requires a
nonzero target, a fixed label count $b$, orthonormal pairs at each window,
unit amplitudes, and the stated error against the normalized positive-part
state. This specifies hypotheses for a conditional theorem; it is not
claimed to be the unique class intended by the source's remark.

Under the additional uniform rate $KN^ke^{-rq}$ on partitions with
$q\le\ell_j\le2q$ at every $1\le q\le N$, the written theorem constructs
$q=\lceil a\log(N/\varepsilon)+B\rceil$, where
$a=\max(k,1)/r$ and $B=\max(\log K,0)/r+L_0+1$. It obtains depth
$C(1+\log(N/\varepsilon))$ with error at most $\varepsilon$, using exact
whole-chain preparation when $q>N$. The uniform rate remains a hypothesis.

The numerical obstruction in
\doclink{TNLean/MPS/Preparation/QualitativeBlockingObstruction} makes the
missing implication precise. Errors $\delta(N,q)=1/(N+1)$ for $q<N$ and
zero for $q\ge N$ decrease with $q$, vanish at whole-ring blocking, and
converge to zero along every schedule. For every fixed $C$, accuracy
$\varepsilon=1/(2(N+1))$ nevertheless requires $q>N-1$ for all sufficiently
large $N$, while $C\log(N/\varepsilon)<N$. These are numerical error
functions, not a realization as optimal MPS errors or a circuit lower bound.

The existing supported-tree theorem assumes a single tensor and a fixed
coordinate virtual support. Neither condition is part of the coherent-pair
approximation condition. In
\doclink{TNLean/MPS/Preparation/NonCoordinateSupport}, the rational virtual
rotation $R=\left(\begin{smallmatrix}3/5&-4/5\\4/5&3/5\end{smallmatrix}\right)$
of the GHZ flags $B^j=R|j\rangle\langle j|$ has $FF^\dagger=I_2$ for its
physical matrix $F$, but admits no injective coordinate support $S$.
Every matrix in its span satisfies $3X_{10}-4X_{00}=0$, which excludes the
matrix unit $E_{00}$ forced by any such $S$. This rules out inferring that
particular support hypothesis from rank bounds alone, not the existence
of a more general tree.

An inhomogeneous tree requires frames and merge maps depending on their
intervals, rather than just their lengths. The generic circuit geometry
already allows position-dependent gates, but the existing homogeneous
supported-tree theorem does not identify its contraction with arbitrary
inhomogeneous polar factors. This is a limitation of that theorem, not a
proved depth obstruction.

A positive local completion step is written in
\doclink{TNLean/MPS/Preparation/IsometricPolarMerge}, using the pinned
QICLean Gram-equality unitary-extension theorem. Every rectangular matrix
$B$ with input dimension at most its output dimension admits $B=XP$ with
$X$ an isometry and $P=(B^\dagger B)^{1/2}$, without injectivity. If $B$
lies in the range of an isometry $J$, the completion can be chosen inside
that range: $B=(JW)P$ with both $W$ and $JW$ isometric. Thus two unrelated
child isometries and an arbitrary compressed parent matrix admit a full
binary isometric merge. This removes the need for a common coordinate
support at that local step.

\section{Algebraic tree continuation}
The interval-dependent construction is written in
\doclink{TNLean/MPS/Preparation/InhomogeneousPolarTree}. If the child
physical matrices factor as $F_A=XP$ and $F_B=YQ$, the actual concatenated
matrix is $(X\otimes Y)C$, where
$C_{(u,v),(a,b)}=\sum_cP_{u,(a,c)}Q_{v,(c,b)}$.
Thus the parent range inclusion follows from the matrix-product
contraction; it is not an extra range hypothesis. Recursive completion
inside that range yields an isometric polar tree on any feasible binary
geometry, with $XP=F$ and $X\Pi=V$ at its root. In particular, if
$s\ge1$ and $d^s\ge D^2$, arbitrary site-dependent chains of length
$s2^h$ have exact polar trees of height $h$ and leaves of at most $s$
sites. No injectivity or common coordinate support is imposed.

This is an algebraic factorization. The generalized theorem in
\doclink{TNLean/MPS/Preparation/IsometryTreePreparation} implements
arbitrary position-dependent unitary node and leaf gates in depth
$C(h+1)$ on disjoint tree layouts. Its outcome-history scalar is
independent of the input on the zero-scratch subspace, hence common
across coherent sectors. The existing theorem for maps depending only
on interval lengths is retained as a specialization.

An end-to-end preparation theorem must still choose unitary extensions
of the actual interval isometries and identify their entire placed gate
product on encoded inputs with the algebraic tree map. This includes
all physical scratch sites and the correspondence between the binary
geometry and the disjoint register layouts. The generic gate compiler
does not by itself prove that coefficient identity.

\section{What a sharp uniform rate would require}
The source's analytic bounds (23) and (S11)--(S16) concern a fixed
translation-invariant tensor and its normal-sector decomposition.
They give admissible exponents $0<\gamma<1/2$, not an optimality theorem
for the approximation exponent. For general non-normal TI tensors, (S16)
also contains an inter-sector term exponential in $N$; (S12) absorbs it
under $q=o(N)$. These length scales and constants belong to a fixed tensor,
not to an arbitrary inhomogeneous sequence. The inhomogeneous random-MPS discussion
in Supplemental Material~3 reports exponential decay per block in
1000 samples, with fitted rates varying between samples, in (S42).
It does not prove one positive rate and one prefactor uniform in length
and in an entire specified inhomogeneous class. Its numerical observations
cannot by themselves select or justify a sharp uniform exponent for the
long-range sequence class left unspecified by the remark.

\doclink{TNLean/MPS/Preparation/ExponentialBlockingThreshold} instead
isolates the exact numerical consequence of a supplied certificate:
$KN^ke^{-rq}\le\varepsilon$ holds precisely when
$q\ge(\log K+k\log N-\log\varepsilon)/r$.
The nonnegative ceiling is the least integral scale satisfying that
certificate. Retaining the source's number-of-blocks factor gives the
exact criterion
$K(N/q)e^{-rq}\le\varepsilon$ if and only if
$rq+\log q\ge\log K+\log N-\log\varepsilon$ for $q>0$.
These exact thresholds are not lower bounds on actual MPS approximation
errors or on preparation depth, and they do not establish a rate for an
unspecified sequence class.

For an unrestricted reading of the source remark, a class of sequences
must still be specified and a uniform rate proved for that class, with a
uniform label count. The conditional theorem establishes how such a rate
gives a logarithmic scale; it does not establish the rate itself. The
doubly-logarithmic claim additionally requires the end-to-end coefficient
identity just described. These obligations are
distinct from the short-range rate work. Issue~\#10573 remains open.

\begin{thebibliography}{9}
\bibitem{Malz2024Preparation}
D.~Malz, G.~Styliaris, Z.-Y.~Wei and J.~I.~Cirac,
\emph{Preparation of matrix product states with log-depth quantum circuits},
Phys. Rev. Lett. \textbf{132}, 040404 (2024),
\url{https://arxiv.org/abs/2307.01696}.
\end{thebibliography}
\end{document}
